% ============================================================
% ch04.tex —— 第 4 章 素数有无穷多个
% ============================================================
\chapter{素数有无穷多个}

先做一次人口普查。拿出 $100$ 以内的素数表（自己筛一遍更有感觉：写下 $2$ 到 $100$，划掉 $2$ 的倍数（保留 $2$），再划掉 $3$ 的、$5$ 的、$7$ 的——剩下的就是素数，这张流程叫\textbf{埃拉托斯特尼筛法}，两千三百年前的希腊人发明的"批量查户口"）：
\[
2,3,5,7,11,13,17,19,23,29,31,37,41,43,47,53,59,61,67,71,73,79,83,89,97 ,
\]
共 $25$ 个。继续数（这活儿可以外包给计算机）：$1000$ 以内 $168$ 个；$10000$ 以内 $1229$ 个。密度在降：前一百个数里四分之一是素数，前一万个数里只剩八分之一。一个自然的怀疑浮上心头：

\begin{kaiwei}
素数是不是越来越稀，稀到最后干脆\textbf{没有了}？会不会存在"世界上最大的素数"，比它更大的数全是合数？
\end{kaiwei}

这个问题在两千三百年前就有答案了。欧几里得的回答是我心目中全人类数学史上最漂亮的证明之一——它只有几行，却示范了一种全新的思维武器。

\section{欧几里得的魔术}

\begin{dingli}[欧几里得，《几何原本》第九卷命题 20]
素数有无穷多个。
\end{dingli}

\begin{proof}
反证法。\textbf{第一步：假设敌人存在。}假设素数不是无穷多——即素数只有有限个，全部一个不漏地列在名单上：
\[
p_1, \; p_2, \; p_3, \; \dots, \; p_n .
\]
（比如 $n = 5$：$2, 3, 5, 7, 11$。"有限"的意思就是能这样全部列完。）

\textbf{第二步：造一个怪物。}把名单上所有素数乘起来，再加 $1$：
\[
N = p_1 \times p_2 \times p_3 \times \cdots \times p_n + 1 .
\]

\textbf{第三步：审问怪物。}$N$ 比名单上任何素数都大，所以它自己不在名单上。那它是什么？由第 3 章，$N > 1$，要么是素数，要么能拆成素数乘积。逐一审查：

\emph{情况一}：$N$ 是素数。那么名单漏了它——"完整名单"不完整，矛盾。

\emph{情况二}：$N$ 是合数。设 $q$ 是它的一个素因数。$q$ 是素数，按假设就该在名单上，也就是 $q$ 等于某个 $p_i$。可是用 $p_i$ 去除 $N$ 会怎样？乘积部分 $p_1\cdots p_n$ 被 $p_i$ 整除（$p_i$ 是乘数之一），剩下加的那个 $1$ 不被整除——所以
\[
N \equiv 1 \pmod{p_i} \quad\text{（第 1 章的余数替换：加 $1$ 之后余数恰好多 $1$）},
\]
$p_i$ 除 $N$ 余 $1$，除不尽！于是 $q$ 不可能等于任何 $p_i$——它是个\textbf{名单外的素数}。名单又漏了，还是矛盾。

\textbf{第四步：矛盾收网。}无论哪种情况，"有限完整名单"都必然有漏。假设崩塌，素数有无穷多个。
\end{proof}

停下来，回味三秒。这个证明最反直觉的地方：\textbf{我们没有造出任何一个新的素数}。看实例——用前六个素数执行构造：
\[
N = 2\times3\times5\times7\times11\times13 + 1 = 30030 + 1 = 30031 = 59 \times 509 ,
\]
$N$ 本身\textbf{不是素数}！但它的素因数 $59$ 和 $509$ 都不在原名单 $\{2,3,5,7,11,13\}$ 里——名单外的素数确实被"逼"出来了，只是藏在 $N$ 的因数里，而不是 $N$ 自己。欧几里得证明的不是"这里有个新素数"，而是一件更抽象也更强大的事：\textbf{任何名单必有遗漏}。"无穷多"的真正含义就是这个——不是"数不完"（数不清的东西多了），而是"任何试图收编全部的努力必败"。

顺带一提：使"$N$ 恰好自己是素数"的名单也很多，比如前四个：$2\times3\times5\times7+1 = 211$（素数）；前五个：$2311$（素数）。运气时好时坏，但\textbf{证明从不依赖运气}——这就是好数学的标志。

\section{反证法：数学里的"归谬"战术}

第一次正式用反证法，值得把说明书读全。它和"直接证明"的根本区别在于进攻方向：

\begin{center}
\begin{tabular}{lp{0.36\textwidth}p{0.36\textwidth}}
\toprule
 & 直接证明 & 反证法 \\
\midrule
假设 & 结论成立 & 结论\textbf{不}成立（敌人当权） \\
推进 & 从已知推向结论 & 从"敌人当权"出发合法推理 \\
终点 & 到达结论 & 撞上荒谬（与已知事实或自身矛盾） \\
判决 & 得证 & 敌人不成立，原命题得证 \\
\bottomrule
\end{tabular}
\end{center}

日常语言里你早就会这招："如果昨晚是你偷吃了蛋糕，冰箱里就该少一块——冰箱完好无损，所以不是你。"数学只是把它磨得更锋利。用它的时机也有规律可循：\textbf{当结论说"不存在/不可能/无穷"时}（要证的东西是个"无"，直接构造没法下手，只好假设"有"然后逼它自爆）。本书后面还会多次请它出场：第 10 章证"零因子没有倒数"、第 12 章判"尺规作图死刑"、第 21 章证"调和级数发散"，招式全同。

\begin{cidian}
土名"假设敌人存在，逼他自曝"$\to$ 正式名\textbf{反证法}（proof by contradiction），拉丁文叫 \textit{reductio ad absurdum}，意为"归结为荒谬"。英国数学家哈代（《一个数学家的辩白》作者）称它为"棋手弃子取势中最锋利的一着"。
\end{cidian}

\section{更强的结论：素数之间可以隔多远？}

"无穷多个"只保证素数永不熄火，但没说它们挤不挤。两个方向的问题都值得问，而且答案都很惊人。

\textbf{方向一：素数能稀到什么程度？}能稀到任意荒谬的程度——想要多长的"素数荒漠"，都能造出来。

\begin{dingli}
对任何 $n$，存在连续 $n-1$ 个整数，全部是合数。
\end{dingli}

\begin{proof}
看这串数：
\[
n! + 2,\quad n! + 3,\quad n! + 4,\quad \dots,\quad n! + n ,
\]
其中 $n! = 1\times2\times3\times\cdots\times n$（阶乘）。任取其中一个 $n! + k$（$2 \le k \le n$）：$n!$ 里含有乘数 $k$，所以 $k \mid n!$，于是
\[
n! + k = k\left(\frac{n!}{k} + 1\right),
\]
两个都大于 $1$ 的整数相乘——合数。要 $1000$ 个连续合数？取 $n = 1001$，那 $999$……准确说 $n!+2$ 到 $n!+n$ 共 $n-1$ 个连续合数，取 $n = 1001$ 即得 $1000$ 个连在一起。荒漠随心定制。
\end{proof}

\textbf{方向二：素数能密到什么程度？}孪生素数——相差 $2$ 的素数对：
\[
(3,5),\ (5,7),\ (11,13),\ (17,19),\ (29,31),\ (41,43),\ (59,61),\ (71,73),\ \dots
\]
（$100$ 以内恰好 $8$ 对，动手找全。）它们会不会无穷多对？——\textbf{没有人知道}。这就是著名的孪生素数猜想，悬而未决超过一百七十年。目前人类最好成绩：2013 年张益唐证明存在无穷多对相差\textbf{不超过七千万}的素数；全球数学家随后协力把差距压到了 $246$。离 $2$ 还差最后 $244$，但按研究惯例，这最后一步可能还要几十年，也可能明天就被踩破。

于是图景变得非常立体：\textbf{素数既永不枯竭，又能稀疏到任意地步；既会成对出现（至少目前看起来没完没了），又会在大片区域绝迹}。那么真正的分布规律是什么？素数与素数之间的"间隔"有没有平均节奏？——这个问题欧几里得答不了，人类又啃了两千年，答案要到第 23 章才揭晓，而且要用上微积分。先把悬念挂起来，记住此刻的心情。

\begin{lishi}
目前已知最大的素数几乎都是\textbf{梅森素数}（形如 $2^p - 1$），因为存在专门针对它们的快速素性检验法（卢卡斯--莱默检验，只需对这类数做一连串"平方再减二"的取模运算）。2018 年 GIMPS（互联网梅森素数大搜索，一个全世界志愿者捐出电脑闲时的项目）发现了第 51 个梅森素数 $2^{82589933} - 1$，超过两千四百万位。对照欧几里得的五行证明，有种奇特的和谐：\textbf{找大素数是工程，证"素数无穷"是数学}——前者的世界纪录会不断刷新，后者两千年纹丝不动。
\end{lishi}

\begin{dongshou}
\begin{itemize}
  \yisi 用前四个素数执行欧几里得构造：$N = 2\times3\times5\times7+1 = 211$，用第 3 章试除法（试到 $\sqrt{211} \approx 14.5$，即 $2,3,5,7,11,13$）验证 $211$ 是素数——名单漏掉了它。再用前五个（$N = 2311$，试到 $\sqrt{2311}\approx 48$）验一次。最后用前六个（$N = 30031$，试到 $\sqrt{30031} \approx 173$）——运气用完了，亲手把它拆成 $59\times509$。
  \yier 用 $7! = 5040$ 造出一段连续合数：写出 $7!+2$ 到 $7!+7$ 这六个数（$5042, 5043, 5044, 5045, 5046, 5047$），并给每一个配上"因数身份证"（如 $5044 = 4\times(5040/4+1) = 4\times1261$）。
  \yier $100$ 以内的孪生素数共 $8$ 对，找全它们。（上文已列出前 $8$ 对，请先遮住上表自己筛一遍再对答案。）
\end{itemize}
\end{dongshou}

\begin{shaonao}
欧几里得的招式可以改装出更细的结论：\textbf{形如 $4k+3$ 的素数有无穷多个}。证明骨架：反设 $4k+3$ 型素数只有 $q_1, \dots, q_m$，构造
\[
M = 4\,q_1q_2\cdots q_m - 1 = 4(q_1\cdots q_m - 1) + 3 ,
\]
它除以 $4$ 余 $3$。审问它：若 $M$ 的素因数全是 $4k+1$ 型，则它们的乘积也是 $4k+1$ 型（关键引理：$(4a+1)(4b+1) = 16ab+4a+4b+1 = 4(\dots)+1$，请展开验证）——与"$M$ 余 $3$"矛盾。所以 $M$ 必有 $4k+3$ 型素因数，而它除不尽 $M$（余数不为零，逐个检验），又不在名单上，矛盾收网。同一招还能证"$4k+1$ 型素数无穷多"（但构造要换成平方数 $M = (2\,q_1\cdots q_m)^2 + 1$，用到"平方数除以 $4$ 余 $0$ 或 $1$"）——这已经摸到了狄利克雷定理（"$ak+b$ 型素数无穷多，只要 $\gcd(a,b)=1$"，1837 年，证明要用到第 22 章的欧拉机器）的门槛。
\end{shaonao}

\begin{chengshi}
严格说，欧几里得原文的写法与本书略有不同：他没有用反证法的形式化语言，而是直接说"任取有限个素数，都能从中逼出一个新素数"（命题 20 的原始表述是"素数的个数多于任何指定的数目"）。本书采用现代反证法的包装，内容等价。另外 $n!+k$ 那段证明里"$k\left(\frac{n!}{k}+1\right)$ 是合数"默认了两个因子都大于 $1$：$k \ge 2$，且 $\frac{n!}{k} + 1 \ge 2$（$n \ge 2$ 时成立），都交代清楚了。
\end{chengshi}

素数的"人口规律"悬念挂起。下一章我们把第 1 章的余数世界升级成一套完整的操作系统——然后请出一位两千年前的中国将军，用这套系统点兵。
